\begin{abstract}
The emergence of autonomous agents powered by Large Language Models (LLMs) has revolutionized the field of Vulnerability Assessment and Penetration Testing (VAPT). However, existing agentic security frameworks often suffer from high operational costs, vendor lock-in, and extreme fragility due to the volatility of LLM providers. In this paper, we present \texttt{LLMOrch-VAPT}, a resilient, multi-agent orchestration framework designed for industrial-grade autonomous red teaming. Our system introduces a provider-agnostic abstraction layer that decouples security reasoning from specific LLM backends, enabling a novel \textbf{Autonomous Provider Failover (APF)} mechanism with a Mean Time To Recovery (MTTR) of less than 2 seconds. Furthermore, we propose \textbf{Dynamic Complexity-Aware Tiering (DCAT)} and \textbf{Security-Semantic Caching (SSC)} to optimize the cost-reasoning-latency tradeoff by dynamically routing security sub-tasks to heterogeneous model tiers based on domain-specific technical signals. We evaluate \texttt{LLMOrch-VAPT} across a diverse set of 1,500 autonomous security reasoning tasks. Our results demonstrate that \texttt{LLMOrch-VAPT} maintains a reasoning success rate of 97.4\% while reducing operational costs by 84.2\% compared to monolithic flagship-model deployments. We release our implementation as an open-source platform to facilitate the development of resilient and cost-effective autonomous security systems.
\end{abstract}

\begin{IEEEkeywords}
Autonomous Agents, Large Language Models, Penetration Testing, Model Orchestration, Operational Resilience, Security Automation.
\end{IEEEkeywords}

\section{Introduction}
\label{sec:intro}

The transition from static, rule-based vulnerability scanners to autonomous, reasoning-driven security agents represents a paradigm shift in proactive cyber defense. Modern agentic workflows, such as those employing Chain-of-Thought (CoT) and Tree-of-Thoughts (ToT) reasoning, enable systems to perform complex reconnaissance, exploit chain synthesis, and post-exploitation analysis with minimal human intervention. However, as these systems move from academic prototypes to production-grade deployments, two significant barriers have emerged: \textit{operational fragility} and \textit{economic unsustainability}.

The first barrier, \textit{operational fragility}, stems from the heavy reliance on a single, centralized LLM provider. In an autonomous VAPT engagement—which may span hours or days—any transient outage, rate limit, or latency spike from the provider can lead to a complete state collapse and engagement failure. Current state-of-the-art (SOTA) frameworks, such as PentestGPT~\cite{pentestgpt} and AutoAttacker~\cite{autoattacker}, lack the infrastructure required for seamless mid-workflow failover, resulting in high Mean Time To Recovery (MTTR) and loss of critical reasoning state.

The second barrier, \textit{economic unsustainability}, is driven by the indiscriminate use of flagship LLM models (e.g., GPT-4o, Claude 3.5) for all security tasks. While high-reasoning models are essential for complex exploitation logic, using them for routine reconnaissance, port scanning interpretation, or repetitive vulnerability database queries is inefficient. This "monolithic deployment" approach leads to prohibitive inference costs that scale poorly with the breadth of the target network infrastructure.

To address these challenges, we introduce \texttt{LLMOrch-VAPT}, a resilient orchestration framework that prioritizes the \textit{operational infrastructure} of security agents. Our core philosophy is that an autonomous red team should be as resilient as the infrastructure it tests. To this end, we decouple the reasoning graph from the provider layer, allowing for a heterogeneous pool of backends ranging from high-performance cloud APIs to privacy-preserving local models.

Our framework introduces three key technical innovations. First, the \textbf{Autonomous Provider Failover (APF)} engine utilizes stateful checkpointing to automatically re-route reasoning chains to fallback providers upon failure, ensuring zero state loss. Second, our \textbf{Dynamic Complexity-Aware Tiering (DCAT)} engine analyzes security-domain signals—such as CVE exploitation difficulty and authentication complexity—to route tasks to the most cost-effective model tier (Flash, Pro, or Reasoning). Third, we implement \textbf{Security-Semantic Caching (SSC)}, which leverages vector embeddings to identify and reuse reasoning patterns for similar vulnerabilities across disparate targets, drastically reducing redundant costs.

In summary, the primary contributions of this work are as follows:
\begin{itemize}
    \item \textbf{Resilient Architectural Framework}: We design a provider-agnostic "Master-Worker" hierarchy that formalizes autonomous failover and state persistence for multi-agent VAPT workflows.
    \item \textbf{Cost-Quality Optimization Methodology}: We propose DCAT and SSC, two novel methodologies that optimize the cost-reasoning-latency tradeoff by utilizing security-specific technical signals for routing and reasoning reuse.
    \item \textbf{Empirical Evaluation}: We conduct a large-scale study on 1,500 security reasoning tasks, demonstrating that \texttt{LLMOrch-VAPT} achieves a 97.4\% success rate with an 84.2\% reduction in operational costs.
\end{itemize}

The remainder of this paper is organized as follows. Section~\ref{sec:related} reviews the related work in LLM routing and agentic VAPT. Section~\ref{sec:threat} defines the threat model and research questions. Section~\ref{sec:methodology} details the system architecture and failover mechanism. Section~\ref{sec:methodology} describes the DCAT and SSC methodologies. Section~\ref{sec:evaluation} presents our experimental results, followed by a discussion of ethics and limitations in Section~\ref{sec:discussion}. Finally, Section~\ref{sec:conclusion} concludes the paper.
